\documentclass[12pt]{article}
\usepackage[utf8]{inputenc}
\usepackage[T1]{fontenc}
\usepackage[margin=1in]{geometry}
\usepackage{setspace}
\usepackage{xcolor}
\usepackage[colorlinks=true, urlcolor=blue!60!black, linkcolor=blue!60!black]{hyperref}
\setlength{\parindent}{0em}
\setlength{\parskip}{0.45em}
\renewcommand{\baselinestretch}{1.0}
\begin{document}
\begin{center}
{\Large \textit{Obergefell v. Hodges}, 576 U.S. 644 (2015)\par}
\vspace{0.3em}
{\normalsize Case Brief --- One-Page Edition\par}
\end{center}

\textbf{Memory Jogger.} The same-sex marriage case: liberty and equality together require states to license and recognize same-sex marriages (\href{https://www.courtlistener.com/opinion/2812209/obergefell-v-hodges/}{576 U.S. 644}).

\textbf{Facts.} Couples challenged marriage bans in four states, including James Obergefell, whom Ohio would not list as his late husband's surviving spouse.

\textbf{Procedural History.} The Sixth Circuit reversed favorable district court rulings, relying on \textit{Baker v. Nelson}.

\textbf{Votes.} Reversed, 5--4: Kennedy for the Court; Roberts, Scalia, Thomas, and Alito each dissented.

\textbf{Holding.} The right to marry is fundamental and extends to same-sex couples; states must license and recognize their marriages, and \textit{Baker} is overruled.

\textbf{Reasoning.} Rights are identified through ``reasoned judgment,'' not \textit{Glucksberg}'s narrow formula; \textit{Loving} protected the ``right to marry in its comprehensive sense.'' Autonomy, the two-person union, children, and social order apply equally to same-sex couples, and liberty and equality work in ``synergy.'' Roberts asked ``who decides''; Scalia decried a ``judicial Putsch''; Thomas called substantive due process a ``dangerous fiction''; Alito found no deep tradition.

\textbf{Examples.} \emph{Same-Side:} birth certificates must list both spouses (\textit{Pavan}). \emph{Opposite-Side:} age and consanguinity limits survive. \emph{Fence-Sitter:} plural marriage.

\textbf{Critique.} Originalists see judicial values displacing history, and \textit{Dobbs} revived doubts; progressives praise ``equal dignity'' but note no scrutiny tier for sexual orientation.

\textbf{Key Quotation.} ``The identification and protection of fundamental rights is an enduring part of the judicial duty to interpret the Constitution. That responsibility, however, `has not been reduced to any formula.' \textit{Poe v. Ullman}, 367 U.S. 497, 542 (1961) (Harlan, J., dissenting). Rather, it requires courts to exercise reasoned judgment in identifying interests of the person so fundamental that the State must accord them its respect.'' --- Justice Kennedy, Opinion of the Court, 576 U.S. at 664.

\end{document}
